\section{Degenerate flag orbits and compact-group orthogonality}
\label{df-section}

The vanishing of a degenerate term between two different compact
types has a useful direct proof.  We give the flag calculation first,
including the case in which the two isotropic subspaces meet.  We then
identify a common compact unitary group and prove the required
orthogonality.  The argument uses no estimate for a lattice mass and
no asymptotic estimate for an orbital integral.  As usual, an orbital
integral below is considered only when it is absolutely convergent.
Its absolute convergence is a separate analytic assertion about the
functions used to construct the theta columns.

\subsection{Flags, characters, and the determinant}

Let $E/F$ be a CM extension, let a bar denote its nontrivial
$F$-automorphism, and let $V$ be an $E$-vector space of dimension
$k+2r+1$.  Suppose that $h$ is a nondegenerate Hermitian form on $V$,
conjugate linear in its first argument, of signature
$(k+r,r+1)$ at every real place of $F$.  The integers $k,r$ are
nonnegative.  An ordered flag with a negative vector consists of
linearly independent vectors $e_1,\ldots,e_r$ and a vector $z$ such
that
\[
 h(e_a,e_b)=0,\qquad h(e_a,z)=0,\qquad h(z,z)=-T,
 \qquad T\in F\text{ totally positive}.
\]
Write
\[
 X_a=\langle e_1,\ldots,e_a\rangle_E,\quad
 X=X_r,\quad X^+=X+Ez,\quad e_{r+1}=z,
 \qquad X_0=0.
\]
The Hermitian space
\begin{equation}\label{df-positive-quotient}
 Q=(X^+)^\perp/X
\end{equation}
is positive definite of dimension $k$ at every real place.  Indeed,
removing $r$ hyperbolic planes gives signature $(k,1)$ on
$X^\perp/X$, and removing the negative line generated by $z$ gives
\eqref{df-positive-quotient}.  In particular,
\begin{equation}\label{df-isotropic-in-perp}
 y\in(X^+)^\perp,\quad h(y,y)=0\quad\Longrightarrow\quad y\in X.
\end{equation}
All assertions about positivity may be checked at any one complex
embedding above a real place when only an algebraic equality is at issue.

Let $N$ be the subgroup of $U(V,h)$ defined by
\[
 (n-1)X_a\subseteq X_{a-1}\ (1\leq a\leq r),
 \qquad (n-1)X^\perp\subseteq X.
\]
It is a connected unipotent algebraic group over $F$.  In a Witt
basis extending the $e_a$, its matrices are triangular with all
diagonal entries equal to one.  Its Lie algebra consists of the
$E$-linear endomorphisms $Y$ satisfying
\begin{equation}\label{df-lie}
 h(Yv,w)+h(v,Yw)=0,\quad
 YX_a\subseteq X_{a-1},\quad YX^\perp\subseteq X.
\end{equation}
If $Y$ satisfies these conditions, let $c_a(Y)$ be the coefficient
of $e_a$ in $Ye_{a+1}$, and put
\begin{equation}\label{df-differential}
 \ell(Y)=\sum_{a=1}^r c_a(Y),\qquad
 d\psi(Y)=\operatorname{Tr}_{E/F}\ell(Y).
\end{equation}
These coefficients are defined because $Ye_{a+1}\in X_a$ for
$a<r$ and $Yz\in X$.  The function $d\psi$ vanishes on Lie
brackets.  For the coefficients with $a<r$, this follows from
$[Y_1,Y_2]X_{a+1}\subseteq X_{a-1}$.  For the last coefficient,
$[Y_1,Y_2]z\in X_{r-1}$.  Thus a fixed additive character of
$F\backslash\mathbb A_F$, composed with $d\psi\circ\log$, defines
a character of $N(\mathbb A_F)$.  Here exponential and logarithm
are finite polynomial expressions in a triangular matrix; the
commutator terms in their multiplication formula disappear under
$d\psi$.  When $r=0$, the group and the character are trivial.

Use primes for a second flag $(e'_1,\ldots,e'_r,z',T')$.  The flags
are called \emph{relevant} when
\begin{equation}\label{df-relevance}
 d\psi(Y)=d\psi'(Y)
 \qquad(Y\in\operatorname{Lie}N\cap\operatorname{Lie}N').
\end{equation}
This is equality of $F$-linear functionals.  It remains an equality
after extension to any completion of $F$.
Put
\[
 I=X\cap X',\qquad V_0=X^++X'^+,
 \quad
 D_a=\det\bigl(h(e_i,e'_j)\bigr)_{1\leq i,j\leq a},
 \qquad D_0=1,
\]
and define
\begin{equation}\label{df-G}
 G=D_{r+1}\overline{D_{r+1}}-TT'D_r\overline{D_r}.
\end{equation}
Conjugating every $D_a$ leaves $G$ unchanged.  Consequently the
opposite convention for which variable of a Hermitian form is linear
does not change the condition $G=0$.

\subsection{The consequence of relevance on a common isotropic vector}

For $0\ne x\in X$, its level is the unique integer $a$ such that
$x\in X_a\setminus X_{a-1}$; write
$x=\alpha e_a+x_0$ with $0\ne\alpha\in E$ and $x_0\in X_{a-1}$.

\begin{lemma}\label{df-successor}
Suppose that the flags are relevant.  If $0\ne x\in I$ has levels
$a,a'$ and leading coefficients $\alpha,\alpha'$ in the two flags,
then
\begin{equation}\label{df-successor-eq}
 \alpha e_{a+1}-\alpha'e'_{a'+1}\in X_a+X'_{a'}.
\end{equation}
\end{lemma}

\begin{proof}
For $v\in(X_a+X'_{a'})^\perp$, consider
\[
 A_{x,v}(y)=h(v,y)x-h(x,y)v.
\]
This endomorphism is skew Hermitian.  If $i\leq a$, it kills $e_i$.
If $i>a$, it sends $e_i$ into $Ex\subseteq X_a\subseteq X_{i-1}$.
It sends $X^\perp$ into $Ex$.  The same statements hold for the
primed flag.  Hence $A_{x,v}$ belongs to both Lie algebras in
\eqref{df-relevance}.  Directly from the coefficients in
\eqref{df-differential},
\[
 \ell(A_{x,v})=\alpha h(v,e_{a+1}),\qquad
 \ell'(A_{x,v})=\alpha'h(v,e'_{a'+1}).
\]
Set $d=\alpha h(v,e_{a+1})-\alpha'h(v,e'_{a'+1})$.
Replacing $v$ by $sv$ shows that
$\operatorname{Tr}_{E/F}(\bar s d)=0$ for every $s\in E$.
Taking $s=d$ gives $2\bar d d=0$, so $d=0$.
Thus $\alpha e_{a+1}-\alpha'e'_{a'+1}$ is orthogonal to
$(X_a+X'_{a'})^\perp$.  Nondegeneracy of $h$ identifies this double
orthogonal complement with $X_a+X'_{a'}$.
\end{proof}

\begin{lemma}\label{df-negative-not-sum}
Neither $z$ nor $z'$ belongs to $X+X'$.  There is consequently at
most one $c\in E$ with $z'-cz\in X+X'$.
\end{lemma}

\begin{proof}
If $z=x+x'$ with $x\in X$ and $x'\in X'$, then
$h(x,x')=h(x,z)=0$.  Both summands are isotropic, so $h(z,z)=0$,
contrary to $T\ne0$.  Interchanging the flags proves the assertion
for $z'$.  Subtracting two expressions $z'-cz\in X+X'$ proves
uniqueness of $c$.
\end{proof}

\begin{lemma}\label{df-common-levels}
If $I\ne0$ and the flags are relevant, there are a unique integer
$0\leq j<r$ and a unique $c\in E^\times$ such that
\begin{align}
 &\dim_E I=r-j,\qquad X=X_j\oplus I,
       \qquad X'=X'_j\oplus I,\label{df-common-splits}\\
 &z'-cz\in X+X',\label{df-z-relation}
\end{align}
and every nonzero vector of $I$ has the same level in both flags,
with leading coefficients related by $\alpha=c\alpha'$.
The levels that occur in $I$ are exactly $j+1,\ldots,r$.
\end{lemma}

\begin{proof}
Start with a common vector of levels $a,a'$.  As long as $a,a'<r$,
write \eqref{df-successor-eq} as
\[
 \alpha e_{a+1}-\alpha'e'_{a'+1}=s+s',
 \qquad s\in X_a,\quad s'\in X'_{a'}.
\]
Then $y=\alpha e_{a+1}-s=\alpha'e'_{a'+1}+s'$ is a common
vector of levels $a+1,a'+1$ with the same two leading coefficients.
If $a\ne a'$, this procedure reaches a stage where just one level
is $r$.  The next use of \eqref{df-successor-eq} puts either $z$ or
$z'$ in $X+X'$, contradicting Lemma~\ref{df-negative-not-sum}.
Thus $a=a'$.  Continuing to level $r$ gives
$z'-(\alpha/\alpha')z\in X+X'$.  Lemma~\ref{df-negative-not-sum}
makes this ratio independent of the original common vector.

Let $j+1$ be the least level occurring in $I$.  The construction
just given produces a vector of $I$ at every larger level.  No
nonzero vector of $I$ lies in $X_j$ or $X'_j$.  Conversely, choose
one vector at each level $j+1,\ldots,r$.  Successively canceling
the leading coefficient of a vector of $X$ shows that these vectors
together with $e_1,\ldots,e_j$ span $X$.  The same reasoning applies
to $X'$.  The direct sums in \eqref{df-common-splits} follow, and
their dimensions give $\dim_E I=r-j$.  This also proves the
asserted uniqueness of $j$.
\end{proof}

\subsection{The complete normal form needed for orthogonality}

\begin{theorem}\label{df-normal-form}
Suppose that the flags are relevant and $G=0$.  Then there are
$0\leq j\leq r$ and $c\in E^\times$ satisfying
\eqref{df-common-splits} and \eqref{df-z-relation}, with
\begin{equation}\label{df-normal-form-values}
 D_j\ne0,\qquad c\bar c T=T'.
\end{equation}
The radical of the restriction of $h$ to $V_0$ is $I$,
\[
 \dim_E V_0=r+j+1,
 \qquad
 \operatorname{sig}(V_0/I)=(j,j+1).
\]
Moreover,
\begin{equation}\label{df-common-perp}
 V_0^\perp\cap X=I=V_0^\perp\cap X',\quad
 V_0^\perp+X=(X^+)^\perp,\quad
 V_0^\perp+X'=(X'^+)^\perp.
\end{equation}
If $I\ne0$, the common-level and leading-coefficient assertions of
Lemma~\ref{df-common-levels} hold.  If $I=0$, one has $j=r$.
\end{theorem}

\begin{proof}
First suppose $I\ne0$, and take $j,c$ from
Lemma~\ref{df-common-levels}.  The pairing between $X_j$ and $X'_j$
is nondegenerate.  Indeed, a vector $x\in X_j$ orthogonal to $X'_j$
is orthogonal to all of $X'$, since $I\subseteq X$.  The relation
$z'-cz\in X+X'$ also makes it orthogonal to $z'$.  It is isotropic,
so \eqref{df-isotropic-in-perp}, applied to the primed flag, gives
$x\in X'$.  Thus $x\in X_j\cap I=0$.  This proves $D_j\ne0$,
including the convention $D_0=1$.

It remains to obtain the same conclusions when $I=0$.
The pairing between $X^+$ and $X'^+$ is then nondegenerate.
For if $x+az\in X^+$ is orthogonal to $X'^+$, its norm is
$-|a|^2T$ at every real place.  On $(X'^+)^\perp$ the form is
positive semidefinite with radical $X'$.  Therefore $a=0$ and
$x\in X'$, whence $x=0$.  Thus $D_{r+1}\ne0$.
Equation~\eqref{df-G} and $G=0$ imply $D_r\ne0$.

The space $X+X'$ is now nondegenerate of signature $(r,r)$.
There are unique $x\in X$, $x'\in X'$ such that
\[
 z_0=z-x\in(X+X')^\perp,\qquad
 z'_0=z'-x'\in(X+X')^\perp.
\]
They are obtained by solving the nondegenerate cross-pairing
systems.  Their norms are $-T$ and $-T'$.  Subtracting the
corresponding linear combinations of the first $r$ rows and columns
in the matrix defining $D_{r+1}$ gives
\[
 D_{r+1}=D_r h(z_0,z'_0).
\]
The space $(X+X')^\perp$ has signature $(k,1)$.  Write
$z'_0=cz_0+w$ with $h(z_0,w)=0$.  Its orthogonal complement to the
negative line $Ez_0$ is positive definite.  Consequently
\[
 |h(z_0,z'_0)|^2=TT'+T h(w,w)
\]
at every real place.  The equality $G=0$ forces $w=0$.
The scalar $c=h(z_0,z'_0)/(-T)$ belongs to $E$ and is nonzero.
This gives \eqref{df-z-relation}, with $j=r$.

We can finish both cases together.  Let
$H_0=X_j+X'_j$, which is a nondegenerate hyperbolic space of
signature $(j,j)$, and note that
$X+X'=H_0\oplus I$ with $I$ orthogonal to $H_0$.
Subtract from $z$ a vector of $X_j$ so that the resulting vector
$z_0$ is orthogonal to $H_0$.  This leaves its norm equal to $-T$;
it is also orthogonal to $I$.  Similarly subtract from $z'$ a
vector of $X'_j$ to obtain $z'_0\perp H_0$, still of norm $-T'$.
The relation \eqref{df-z-relation} implies
$z'_0-cz_0\in H_0+I$.  Since this difference is orthogonal to the
nondegenerate space $H_0$, it lies in $I$.
Comparing its norm with $-T'$ proves $c\bar c T=T'$.
We have the orthogonal decomposition, with the last summand radical,
\[
 V_0=H_0\mathbin{\perp}Ez_0\mathbin{\perp}I.
\]
It proves the assertions about $V_0$.  In particular
$V_0^\perp\cap X=I$ and similarly for $X'$.
The inclusion $V_0^\perp+X\subseteq(X^+)^\perp$ is immediate.
Both sides have dimension $k+r$, since
\[
 \dim(V_0^\perp+X)
 =(k+r-j)+r-(r-j)=k+r.
\]
This proves the other two equalities in \eqref{df-common-perp}.
The argument includes $r=0$, in which case it is just the equality
case of the reverse Cauchy--Schwarz inequality in signature $(k,1)$.
\end{proof}

\begin{corollary}\label{df-common-compact}
Under the hypotheses of Theorem~\ref{df-normal-form}, there is a
positive definite $E$-subspace $W\subseteq V_0^\perp$ of dimension
$k$ such that
\[
 V_0^\perp=W\oplus I.
\]
The maps induced by inclusion are isometries
\begin{equation}\label{df-quotient-isometries}
 W\simeq (X^+)^\perp/X,
 \qquad W\simeq (X'^+)^\perp/X'.
\end{equation}
The subgroup $U(W)$, acting as the identity on $W^\perp$, fixes
$V_0$ pointwise.  At infinity it is a compact group
$K_W\simeq\prod_{v\mid\infty}U(k)$, acting on the two quotients
in \eqref{df-quotient-isometries} by the same standard unitary
action, up to conjugation by a unitary change of basis.
\end{corollary}

\begin{proof}
The radical of $V_0^\perp$ is $I$.  In $I^\perp/I$ the form has
signature $(k+j,j+1)$, while the nondegenerate subspace $V_0/I$
has signature $(j,j+1)$.  Hence its orthogonal complement
$V_0^\perp/I$ is positive definite of dimension $k$.
Choose an $E$-linear complement $W$ of $I$ in $V_0^\perp$.
The form on $W$ is positive definite at every real place.
The injectivity and surjectivity of the maps in
\eqref{df-quotient-isometries} follow from
\eqref{df-common-perp}.  They preserve the Hermitian forms because
the subspaces being divided out are radicals.  Finally,
$V=W\perp W^\perp$, and $V_0\subseteq W^\perp$.
\end{proof}

\subsection{Averaging over the unipotent quotient}

We spell out a measure issue in applying the common compact group.
Its elements need not preserve the particular integral lattices
that define an arithmetic subgroup.  The appropriate operation is
first to integrate over the quotient of the \emph{real unipotent}
group by its arithmetic subgroup.

\begin{lemma}\label{df-unipotent-lattice}
Let $R$ be a unipotent algebraic group over a number field $F$ and
let $K_f\subseteq R(\mathbb A_{F,f})$ be a compact open subgroup.  Then
\[
 \Gamma=R(F)\cap K_f
\]
is a discrete cocompact subgroup of $R(F\otimes_{\mathbb Q}\mathbb R)$.
The group $R(F\otimes_{\mathbb Q}\mathbb R)$ is unimodular.
\end{lemma}

\begin{proof}
In characteristic zero, the exponential and logarithm of a
unipotent matrix are finite polynomials, inverse to one another.
The Lie algebra has a central series defined over $F$, with
additive $F$-vector spaces as successive quotients.  Choose an
$F$-basis compatible with a refinement of this series.  Ordered
exponential coordinates give polynomial coordinates on $R$; in
these coordinates multiplication is triangular, with the new
coordinate equal to the sum of the two coordinates plus a
polynomial in earlier coordinates.

Clearing the finitely many denominators in these polynomial
expressions, and shrinking the resulting fractional ideals if
necessary at the finitely many places where $K_f$ is prescribed,
gives coordinate lattices contained in $\Gamma$.  More explicitly,
each one-parameter subgroup $\exp(tY)$ contains
$\exp(\mathfrak a Y)\subseteq\Gamma$ for some nonzero fractional
ideal $\mathfrak a$; this follows by inserting $t\in\mathfrak a$
in the finite polynomial for $\exp(tY)$ and requiring membership
in $K_f$.  Successive multiplication by these elements reduces
each coordinate to a bounded fundamental parallelepiped for the
image of its ideal in $F\otimes\mathbb R$.  The triangular form
ensures that coordinates already reduced are unchanged.  The
resulting bounded closed coordinate box is compact and meets every
left $\Gamma$-coset.  Discreteness follows because $\log K_f$ is
contained in a product of fractional-ideal completions and the
diagonal image of each fractional ideal is discrete at infinity.
Finally, the adjoint action of a unipotent element on the Lie
algebra is unipotent, so its determinant is one.  This proves
unimodularity.
\end{proof}

Let $G_\infty=U(V)(F\otimes\mathbb R)$, and let
$R_\infty=N_\infty\cap N'_\infty$.  This is the real group of a
unipotent $F$-group.  Let $\Gamma\subset R_\infty$ be an arithmetic
subgroup obtained by imposing stabilizers of two integral lattices.
It has the form in Lemma~\ref{df-unipotent-lattice}, after taking
their compact open intersection.  Choose Haar measures $dg,dr$ and
the quotient measure $d\dot g$ characterized by
\[
 \int_{G_\infty}f(g)\,dg
 =\int_{R_\infty\backslash G_\infty}
       \int_{R_\infty}f(rg)\,dr\,d\dot g.
\]
Set $v_R=\operatorname{vol}(\Gamma\backslash R_\infty)$.

\begin{lemma}\label{df-average}
Suppose that a finite-dimensional vector-valued function $B$ on
$G_\infty$ satisfies $B(rg)=\chi(r)B(g)$, where $\chi$ is a
continuous unitary character of $R_\infty$, trivial on $\Gamma$,
and that $B$ is absolutely integrable on
$\Gamma\backslash G_\infty$.
If $\chi\ne1$, its integral is zero.  If $\chi=1$, then
\begin{equation}\label{df-disintegration}
 \int_{\Gamma\backslash G_\infty}B(g)\,dg
 =v_R\int_{R_\infty\backslash G_\infty}B(g)\,d\dot g.
\end{equation}
If a compact subgroup $K\subseteq G_\infty$ normalizes $R_\infty$,
then the quotient measure on $R_\infty\backslash G_\infty$ is
invariant under left multiplication by $K$.
\end{lemma}

\begin{proof}
Apply the displayed integration formula first to scalar absolute
values, and then to each coordinate of $B$, using a measurable
fundamental domain for $\Gamma$ in $R_\infty$.  This gives the
inner integral over $\Gamma\backslash R_\infty$.
If $\chi\ne1$, right translation by an element on which $\chi$
is nontrivial multiplies its character integral by a scalar other
than one.  The integral is therefore zero.  If $\chi=1$, the
inner integral is $v_R B(g)$, proving
\eqref{df-disintegration}.  Absolute convergence justifies these
uses of Fubini.

Conjugation by $k\in K$ multiplies $dr$ by a positive real scalar.
That scalar is a continuous homomorphism from the compact group
$K$ to $\mathbb R_{>0}$, hence is one.  Left multiplication by
$k$ preserves $dg$ as well.  Uniqueness of the quotient measure
in the integration formula proves the last assertion.
\end{proof}

\subsection{Vanishing of the actual compact-type orbital operator}

At each real place $v$ let $H_v,H'_v$ be finite-dimensional
Hilbert spaces carrying representations $\pi_v,\pi'_v$ of the
corresponding compact groups $U(k)$.  Put
\[
 H=\bigotimes_{v\mid\infty}H_v,\qquad
 H'=\bigotimes_{v\mid\infty}H'_v,
\]
with the external tensor-product actions
$\pi((k_v)_v)=\bigotimes_v\pi_v(k_v)$ and
$\pi'((k_v)_v)=\bigotimes_v\pi'_v(k_v)$.
Let
$\Phi:G_\infty\to H$, $\Phi':G_\infty\to H'$ be the functions
attached to the two flags.  A translate of the second function is
already included in the primed notation.  Suppose that
\[
 \Phi(ng)=\psi(n)\Phi(g),\qquad
 \Phi'(n'g)=\psi'(n')\Phi'(g).
\]
The representations of the compact stabilizers on $H,H'$ factor
through the positive quotients \eqref{df-positive-quotient}.
In particular, for the subgroup in Corollary~\ref{df-common-compact},
\[
 \Phi(kg)=\pi(k)\Phi(g),\qquad
 \Phi'(kg)=\pi'(k)\Phi'(g)\qquad(k\in K_W).
\]
For $v\in H$, $w\in H'$, write
$[v,w](\xi)=\langle w,\xi\rangle v$.  With inner products conjugate
linear in the first variable, define the orbital operator
\begin{equation}\label{df-orbital-operator}
 J=\int_{\Gamma\backslash G_\infty}
          [\Phi(g),\Phi'(g)]\,dg.
\end{equation}
The finite-place admissibility conditions on a pair of theta columns
make $\psi\overline{\psi'}$ trivial on $\Gamma$.  This is also the
precise condition needed for the integrand in
\eqref{df-orbital-operator} to descend to the quotient: rational
additive characters have product one, and their two finite-place
factors agree on the common integral stabilizer.  No normalization
of $v_R$ is needed for the following vanishing assertion.

\begin{theorem}\label{df-operator-vanishing}
Suppose that $G=0$ and that \eqref{df-orbital-operator} converges
absolutely.  At every real place let $\pi_v,\pi'_v$
be irreducible representations of $U(k)$, and suppose that they
are nonisomorphic at at least one place.  Then $J=0$.
\end{theorem}

\begin{proof}
If the flags are not relevant, the differential of
$\chi=\psi\overline{\psi'}$ is nonzero on the common real
unipotent group.  To see this globally, a nonzero $F$-valued
differential remains nonzero after applying
$\operatorname{Tr}_{F/\mathbb Q}(a\,\cdot)$ for some $a\in F$,
by nondegeneracy of the trace pairing.  Thus the character at
infinity is nontrivial.  Lemma~\ref{df-average} gives $J=0$.

Suppose the flags are relevant.  Equality of the differentials
implies equality of the characters on the common real unipotent
group, since exponential coordinates cover that connected group.
Use Theorem~\ref{df-normal-form} and
Corollary~\ref{df-common-compact} to choose $K_W$.
It fixes every vector $e_a,e'_a,z,z'$, so it normalizes both
unipotent groups and preserves their characters.  In particular,
it normalizes $R_\infty$.
By \eqref{df-disintegration} and the last part of
Lemma~\ref{df-average}, change of variables $g\mapsto kg$ gives
\[
 J=\pi(k)J\pi'(k)^{-1},\qquad k\in K_W.
\]
Here the conjugation on the second factor is essential:
$[Av,Bw]=A[v,w]B^{-1}$ for unitary $A,B$.
The isometries \eqref{df-quotient-isometries} identify the actions
with the two external tensor products of the specified $U(k)$
representations.  Changes of orthonormal basis only conjugate
one action by an inner automorphism of $U(k)$ and do not change
its isomorphism class.

For completeness, choose a place where the two irreducible
representations are nonisomorphic.  After fixing vectors and
covectors in all other tensor factors, the displayed intertwining
identity turns every matrix entry of $J$ in those factors into an
intertwiner between the two representations at that place.
Such an intertwiner is zero: its kernel and image are invariant,
so a nonzero intertwiner of irreducible representations would be
an isomorphism.  All matrix entries of $J$ vanish, proving $J=0$.
\end{proof}

\begin{corollary}\label{df-points}
Consider two distinct words $u,u'$ in the theta family, with their
specified dominant compact Harish--Chandra parameters
$\lambda_u(v),\lambda_{u'}(v)$.  Suppose these parameters differ
at one real place.  Every absolutely convergent term with $G=0$
in the operator from the $u'$-column space to the $u$-column space
is zero.  This includes the identity term whenever it occurs in
that off-diagonal expansion.
\end{corollary}

\begin{proof}
Number the diagonal coordinates from $0$ to $k-1$ and put
$\rho_a=(k-1)/2-a$.  The specified compact type has a nonzero
vector of diagonal-torus weight $\mu=\lambda-\rho$, and every
one of its weights $\nu$ satisfies
\[
 \mu-\nu=\sum_{a<b}c_{ab}(\varepsilon_a-\varepsilon_b),
 \qquad c_{ab}\in\mathbb Z_{\ge0},
\]
where $\varepsilon_a$ are the coordinate vectors.  This is the
highest-weight condition used to specify the compact types;
it suffices here without a representation-classification theorem.
Indeed, an isomorphism of the two representations would preserve
all diagonal-torus weights.  Hence both $\mu-\mu'$ and
$\mu'-\mu$ would be nonnegative combinations of the displayed
positive roots.  The linear functional
\[
             \ell(x)=\sum_{a=0}^{k-1}(k-a)x_a
\]
takes the strictly positive value $b-a$ on
$\varepsilon_a-\varepsilon_b$ for $a<b$.  Applying it to
these opposite combinations shows that every coefficient is
zero.  Thus $\mu=\mu'$, and consequently $\lambda=\lambda'$.
Different specified parameters therefore give nonisomorphic
representations.  Apply
Theorem~\ref{df-operator-vanishing}.  The corresponding term in
the theta operator is assembled from $J$ by finite-place scalar
factors and finite linear maps, all of which send the zero operator
to zero.
\end{proof}

In the source, the distinct-word hypothesis supplies different
compact parameters at every real place, so it is stronger than
the hypothesis of Corollary~\ref{df-points}.  The notation for the
source's determinants is the conjugate of one of the conventions
used in the intermediate normal-form calculation; as observed
after \eqref{df-G}, this has no effect on degeneracy or on
\eqref{df-normal-form-values}.  The calculation proves the
geometric and orthogonality content of the assertion called the
points lemma in the source~\cite{alpo-uniform-source}.
It does not prove the diagonal estimate for the sum of the
nonidentity degenerate terms.  That estimate additionally needs
the sizes of the orbital amplitudes, the count of their lattice
data, and the finite-level normalization.


